% Clase 24 --- Onda viajera y el signo del argumento (§3.2).
% El docente dibuja el pulso en t=0 y "el mismo dibujo corrido" en t. La figura
% tiene que dejar ver por que el MENOS mueve la onda hacia la DERECHA: el maximo
% vive donde el argumento se anula.
\begin{center}
\begin{tikzpicture}[scale=1]

  \draw[figaxis] (-0.5,0) -- (9.2,0) node[figlblaux, anchor=north east] {$x$};
  \draw[figaxis] (0,-0.3) -- (0,2.3) node[figlblaux, anchor=south east] {$y$};

  % perfil en t = 0
  \draw[figgray, line width=1pt, smooth, samples=80, domain=-0.2:3.4]
    plot (\x, {1.7*exp(-((\x-1.6)^2)/0.55)});
  \node[figlblaux, anchor=south] at (1.6,1.78) {$t = 0$};
  \filldraw[figgray] (1.6,1.7) circle (1.6pt);

  % el mismo perfil, corrido
  \draw[figline, line width=1pt, smooth, samples=80, domain=3.4:8.9]
    plot (\x, {1.7*exp(-((\x-6.5)^2)/0.55)});
  \node[figlbl, figblue, anchor=south] at (6.5,1.78) {$t > 0$};
  \filldraw[figblue] (6.5,1.7) circle (1.6pt);

  % el corrimiento
  \draw[figgray, dashed, line width=0.6pt] (1.6,0.12) -- (1.6,1.68);
  \draw[figgray, dashed, line width=0.6pt] (6.5,0.12) -- (6.5,1.68);
  \draw[figred, line width=0.9pt,
        {Latex[length=2mm,width=1.5mm]}-{Latex[length=2mm,width=1.5mm]}]
    (1.6,0.42) -- (6.5,0.42);
  \node[figlbl, figred, anchor=south] at (4.05,0.48) {$\Delta x = v\,t$};

  \node[figlbl, anchor=north] at (1.6,-0.12) {$y(0,x) = f(x)$};
  \node[figlbl, anchor=north] at (6.5,-0.12) {$y(t,x) = f(x - vt)$};

\end{tikzpicture}\\[0.5em]
{\small\itshape El signo menos desconcierta a todo el mundo: parece que para
correr la onda hacia la derecha habría que \emph{sumar}. La verificación se ve
mejor siguiendo un punto de la forma que sea fácil de rastrear, por ejemplo el
máximo. Si se elige el origen de modo que en $t=0$ el máximo esté en $x = 0$,
entonces en cualquier instante posterior el máximo sigue estando donde el
argumento de $f$ se anula, es decir en $x - vt = 0$, o sea en $x = vt$: cada vez
más a la derecha. Con $f(x + vt)$ el mismo razonamiento da $x = -vt$ y la onda
viaja hacia la izquierda. Lo que hace \enquote{viajera} a esta onda es que la
\textbf{forma no cambia}: el segundo dibujo es una copia calcada del primero,
apenas trasladada. En una cuerda real puede haber una onda hacia la derecha
superpuesta con otra hacia la izquierda, y ahí el perfil sí cambia de aspecto
con el tiempo.}
\end{center}
